\documentclass[11pt]{article}

\usepackage[margin=1in]{geometry}
\usepackage{amsmath,amssymb,amsthm}
\usepackage{mathtools}
\usepackage{microtype}
\usepackage{hyperref}
\usepackage[numbers,sort&compress]{natbib}

\newtheorem{theorem}{Theorem}
\newtheorem{proposition}{Proposition}
\newtheorem{lemma}{Lemma}
\newtheorem{definition}{Definition}
\newtheorem{example}{Example}

\title{Constraints on Information Localization and Recoverability\\
in Quantum-Gravitational Bounce Models}
\author{Steven Owens}
\date{Draft scaffold -- not for submission}

\begin{document}
\maketitle

\begin{abstract}
Black-to-white-hole transition and remnant models motivate a precise question: after a quantum-gravitational transition, where is information about the infalling state recoverable? Let a reference system $R$ purify an input $X$, and let an isometry $V:X\rightarrow A\otimes B$ represent a transition into two output sectors. For the resulting pure state on $RAB$, the exact identity $I(R:A)+I(R:B)=2S(R)$ conserves total reference correlations but does not enforce equal partitioning. We therefore replace a universal no-filtering claim with the conditional problem of deriving recoverability bounds under explicit assumptions concerning accessibility, conservation laws, causal structure, remnant state spaces, semiclassical validity, and holographic code-subspace reconstruction. This manuscript develops the channel-level identities and counterexamples, specifies the conditions required for gravitational embedding, and defines a feasibility gate for any observational application.
\end{abstract}

\section{Introduction}
\label{sec:introduction}

State the physical motivation without claiming that unitarity fixes the distribution of information across arbitrary output subsystems. Distinguish global purity, subsystem entropy, and operational recoverability \citep{Page1993Information,HaydenPreskill2007}.

\section{Channel framework}
\label{sec:channel}

\begin{definition}[Reference-assisted transition]
Let $\rho_{RX}$ be a pure state and let $V:\mathcal H_X\rightarrow\mathcal H_A\otimes\mathcal H_B$ be an isometry. Define
\[
\rho_{RAB}=(\operatorname{id}_R\otimes V)\rho_{RX}(\operatorname{id}_R\otimes V^\dagger).
\]
\end{definition}

\begin{proposition}[Mutual-information sum]
For pure $\rho_{RAB}$,
\[
I(R:A)+I(R:B)=2S(R).
\]
\end{proposition}

\begin{proof}
Using purity, $S(RA)=S(B)$ and $S(RB)=S(A)$. Therefore
\begin{align*}
I(R:A)+I(R:B)
&=2S(R)+S(A)+S(B)-S(RA)-S(RB)\\
&=2S(R).
\end{align*}
\end{proof}

\begin{example}[Maximally biased isometry]
For $V|\psi\rangle_X=|\psi\rangle_A\otimes|0\rangle_B$, one has
\[
I(R:A)=2S(R),\qquad I(R:B)=0.
\]
Thus unitarity alone does not imply equal information localization.
\end{example}

\section{Recoverability and decoupling}
\label{sec:recovery}

Define operational recovery tasks, norm conventions, coherent information, and one-shot decoupling criteria \citep{DupuisEtAl2010,FawziRenner2014}. State all assumptions before introducing a conditional theorem.

\section{Operator-algebra and holographic benchmarks}
\label{sec:holography}

Treat reconstruction from specified observable algebras using operator-algebra quantum error correction \citep{BenyKempfKribs2007,AlmheiriDongHarlow2014,Harlow2016}. Use island and entanglement-wedge results only within the controlled models where their assumptions hold \citep{Penington2019,AlmheiriEtAl2019Bulk,AlmheiriEtAl2019Replica}.

\section{Black-to-white-hole model cards}
\label{sec:models}

For each selected model, record the geometry, asymptotics, state space, transition law, radiation sector, remnant sector, and domain of validity. Candidate literature includes black-to-white-hole tunneling and white-hole remnant constructions \citep{HaggardRovelli2014,BianchiEtAl2018}.

\section{Phenomenology feasibility gate}
\label{sec:phenomenology}

Require an explicit spectrum, duration, energy budget, event-rate prescription, source population, propagation model, detector response, background model, and statistical plan. Do not equate $2GM/c^3$ with an observed transient duration without a model-derived connection.

\section{Discussion and limitations}
\label{sec:discussion}

Separate universal channel identities from conditional gravitational results. Document counterexamples, model dependence, inaccessible algebras, and unresolved microscopic dynamics.

\section{Conclusion}
\label{sec:conclusion}

Summarize only results established in the completed sections.

\bibliographystyle{unsrtnat}
\bibliography{../references/references}

\end{document}
